\chapter{المنحنيات}\label{ch:b2:curves}

تصير الهندسة الآن تفاضلية. فالمنحنى نقطة تتحرك في الفضاء؛ ويعطينا
حساب \omterm{def:b2:diffcalc:differential}{التفاضل} سرعتها وتسارعها، وتسأل الهندسة عما هو جوهري --- أي
مستقل عن سرعة قطعنا للمسار. والجوابان هما \emph{\omterm{def:b2:curves:length}{طول القوس}}، الذي
يقيس المسار نفسه، و\emph{\omterm{thm:b2:curves:frenet2d}{الانحناء}}، الذي يقيس مقدار انثنائه. وفي
البعد $3$ يقيس ثابت ثان، هو \emph{\omterm{thm:b2:curves:frenet3d}{الالتواء}}، مقدار خروج المنحنى
عن مستويه بالتفتّل. وأداة مسك الحسابات في هذا كله هي \emph{معلم
فرينيه} المتحرك.

\section{الأقواس الموسّمة}

\begin{definition}[القوس الموسّم]\label{def:b2:curves:arc}
\emph{القوس الموسّم}\index{قوس موسم} من الصنف $\mathcal{C}^k$ ($k \geq 1$) هو تطبيق
$\gamma \colon I \to \R^n$ من الصنف $\mathcal{C}^k$ على فترة $I$. وتكون النقطة $\gamma(t)$
\emph{نظامية} إذا كان $\gamma'(t) \neq 0$، ويكون القوس \emph{نظاميا} إذا كانت
جميع نقطه كذلك. والمستقيم المار بالنقطة $\gamma(t)$ والموجه بالمتجهة
$\gamma'(t)$ هو \emph{المماس}\index{مماس} في نقطة نظامية.
\end{definition}

\begin{definition}[تغيير الوسيط]\label{def:b2:curves:reparam}
\emph{تغيير الوسيط}\index{تغيير الوسيط} من الصنف $\mathcal{C}^k$ هو تقابل
تفاضلي $\theta \colon J \to I$ من الصنف $\mathcal{C}^k$ بين فترتين (مع $\theta' \neq 0$ في كل نقطة).
ويقال إن القوسين $\gamma$ و$\gamma \circ \theta$ \emph{متكافئان}؛ و\emph{القوس
الهندسي}\index{قوس هندسي} (أو المنحنى) هو صف تكافؤ. والمفاهيم
الثابتة عند تغيير الوسيط --- المسار \omterm{def:b2:curves:arc}{والمماس} \omterm{def:b2:curves:length}{وطول القوس} \omterm{thm:b2:curves:frenet2d}{والانحناء} ---
تسمى \emph{هندسية}.
\end{definition}

\begin{remark}
لا يحدد المسار وحده \omterm{def:b2:curves:reparam}{القوس الهندسي}: فالتوسيمان $t \mapsto (\cos t, \sin t)$ على $[0, 2\pi]$ وعلى
$[0, 4\pi]$ لهما الصورة نفسها لكن أحدهما يقطع الدائرة مرة والآخر مرتين.
\omterm{def:b2:curves:reparam}{فالقوس الهندسي} يتذكر تعددية القطع واتجاهه، لا سرعته.
\end{remark}

\begin{example}[السرعة لا تغير شيئا هندسيا]\label{ex:b2:curves:speedinvariance}
لنوسّم دائرة الوحدة بالعلاقة
\[
\gamma(t) = (\cos t^2,\ \sin t^2),
\qquad t \in \intcc0{\sqrt{2\pi}} .
\]
تنمو السرعة $\norm{\gamma'(t)} = 2t$ خطيا، ومع ذلك
\[
L = \int_0^{\sqrt{2\pi}}2t\,\dd t = 2\pi ,
\]
أي الطول نفسه الذي نجده بسرعة ثابتة --- كما تَعِد \cref{thm:b2:curves:lengthinv}، عبر \omterm{def:b2:curves:reparam}{تغيير
الوسيط} $t \mapsto t^2$. \omterm{def:b2:curves:arc}{والمماس}، \omterm{thm:b2:curves:frenet2d}{والانحناء} المحسوب من \cref{prop:b2:curves:kappaformula}، وكل كمية هندسية
أخرى، تتوافق أيضا؛ ولا يستحق النظر إلا $t = 0$، حيث يجعل $\gamma'(0) = 0$
\emph{هذا التوسيم} غير نظامي مع أن المسار دائرة تامة. فالنصوص
الهندسية تتحمل التوسيمات السيئة تحملا سيئا: أعد التوسيم أولا، ثم
استنتج ثانيا.
\end{example}

\begin{example}\label{ex:b2:curves:cusp}
القوس $\gamma(t) = (t^2, t^3)$ من الصنف $\mathcal{C}^\infty$ لكنه ليس نظاميا: $\gamma'(0) = (0, 0)$. ولمساره، وهو
القطع المكافئ نصف التكعيبي $y^2 = x^3$، \emph{نقطة رجوع} في المبدأ:
فملاسة التوسيم لا تمنع شذوذا هندسيا حيث تنعدم السرعة. ولهذا فإن
فرضية النظامية $\gamma' \neq 0$ ليست تجميلية.
\end{example}

\section{طول القوس}

\begin{definition}[طول القوس]\label{def:b2:curves:length}
ليكن $\gamma \colon [a, b] \to \R^n$ قوسا من الصنف $\mathcal{C}^1$. \emph{طوله}\index{طول القوس} هو
\[
L(\gamma) = \int_a^b \norm{\gamma'(t)}\, \dd t ,
\]
حيث $\norm{\cdot}$ هو \omterm{def:b2:nvs:norm}{المعيار} الإقليدي. و\emph{دالة طول القوس}\index{دالة طول القوس}
المرتكزة في $t_0$ هي $s(t) = \int_{t_0}^t
\norm{\gamma'(u)}\,\dd u$.
\end{definition}

\begin{theorem}[الطول هندسي؛ التوصيف بالمضلعات]\label{thm:b2:curves:lengthinv}
\leavevmode
\begin{enumerate}
\item إذا كان $\theta \colon [c, d] \to [a, b]$ تغيير وسيط من الصنف $\mathcal{C}^1$، فإن $L(\gamma \circ \theta) = L(\gamma)$.
\item $L(\gamma)$ هو الحد الأعلى لأطوال المضلعات المرسومة داخله:
\[
L(\gamma)
= \sup\Bigl\{\, \sum_{i=1}^{m} \norm{\gamma(t_i) -
\gamma(t_{i-1})} \;:\; a = t_0 < t_1 < \dots < t_m = b \,\Bigr\}.
\]
\end{enumerate}
\end{theorem}

\begin{proof}
\emph{1.} بتغيير المتغير $t = \theta(u)$ (مجلد السنة الأولى، وهو صالح لأن
$\theta$ رتيبة من الصنف $\mathcal{C}^1$)،
\[
\int_c^d \norm{(\gamma\circ\theta)'(u)}\,\dd u
= \int_c^d \norm{\gamma'(\theta(u))}\,\abs{\theta'(u)}\,\dd u
= \int_a^b \norm{\gamma'(t)}\,\dd t ,
\]
حيث استعملنا $(\gamma\circ\theta)' = \theta'\cdot
(\gamma'\circ\theta)$، وإذا كانت $\theta$ متناقصة فإن إشارة $\theta'$ تمتصها
\omterm{def:b2:structures:sn}{مبادلة} الحدين.

\emph{2.} من أجل أي تجزئة، $\gamma(t_i) - \gamma(t_{i-1}) =
\int_{t_{i-1}}^{t_i}\gamma'(t)\,\dd t$، ومنه فبمتراجحة المثلث في التكاملات
$\norm{\gamma(t_i) - \gamma(t_{i-1})}
\leq \int_{t_{i-1}}^{t_i}\norm{\gamma'}$: فكل مضلع أقصر من $L(\gamma)$، ومنه $\sup \leq L(\gamma)$.

ولإثبات المتراجحة العكسية، ليكن $\varepsilon > 0$. بما أن $\gamma'$ متصلة على
المتراص $[a,b]$، فهي متصلة بانتظام: فيوجد $\delta > 0$ بحيث يكون $\norm{\gamma'(t) - \gamma'(u)} \leq
\varepsilon$
كلما كان $\abs{t - u} \leq \delta$. نأخذ تجزئة خطوتها $\leq \delta$. وعلى كل قطعة، من أجل
$t \in [t_{i-1}, t_i]$،
\[
\gamma(t_i) - \gamma(t_{i-1})
= \int_{t_{i-1}}^{t_i}\gamma'(t)\,\dd t
= (t_i - t_{i-1})\,\gamma'(t_{i-1}) + R_i,
\qquad
\norm{R_i} \leq \varepsilon\,(t_i - t_{i-1}),
\]
لأن $R_i = \int_{t_{i-1}}^{t_i}(\gamma'(t) -
\gamma'(t_{i-1}))\,\dd t$. ومنه
\[
\norm{\gamma(t_i) - \gamma(t_{i-1})}
\geq (t_i - t_{i-1})\norm{\gamma'(t_{i-1})}
- \varepsilon (t_i - t_{i-1}) .
\]
وبالجمع، وبمقارنة $\sum (t_i -
t_{i-1})\norm{\gamma'(t_{i-1})}$ مع $\int_a^b\norm{\gamma'}$ (وهو مجموع ريمان للدالة المتصلة
$\norm{\gamma'}$، ويبعد عن التكامل بأقل من $\varepsilon(b - a)$ من أجل $\delta$ صغيرة بما يكفي،
بالاتصال المنتظم مرة أخرى)، نحصل على مضلع طوله $\geq
L(\gamma) - 2\varepsilon(b - a)$. وبجعل
$\varepsilon \to 0$ يثبت المطلوب.
\end{proof}

\begin{example}[أرخميدس والمضلعات الداخلية]\label{ex:b2:curves:archimedes}
من أجل دائرة الوحدة، يكون طول المضلع المنتظم الداخلي ذي $n$ ضلعا
هو $L_n = 2n\sin\frac\pi n$، ويعطي النشر $\sin x = x -
\frac{x^3}6 + O(x^5)$
\[
L_n = 2\pi - \frac{\pi^3}{3n^2} + O\Bigl(\frac1{n^4}\Bigr) :
\]
فأطوال المضلعات في \cref{thm:b2:curves:lengthinv} تتقارب \emph{تقاربا تربيعيا}. وعدديا:
$L_6 = 6$ (المسدس، وهو يعطي القيمة الفجة $\pi > 3$)، بينما $L_{96} =
192\sin\frac{\pi}{96} \approx 6.28206$ مقابل
$2\pi \approx
6.28319$ --- والخطأ $0.00113$ يوافق القيمة المتنبأ بها $\pi^3/(3\cdot96^2) \approx 0.00112$. ولهذا استطاع
أرخميدس، بمضاعفة المسدس خمس مرات إلى $96$ ضلعا، أن يحصر $\pi$
بثلاثة أرقام يدويا: فكل مضاعفة تقسم الخطأ على أربعة. فالحد الأعلى في
التوصيف بالمضلعات ليس مبلوغا في النهاية فحسب؛ بل هو مبلوغ
\emph{بسرعة}، لأن المنحنى الأملس لا ينفصل عن أوتاره إلا من الرتبة
الثانية.
\end{example}

\begin{theorem}[التوسيم بطول القوس]\label{thm:b2:curves:arclength}
ليكن $\gamma \colon I \to \R^n$ قوسا نظاميا من الصنف $\mathcal{C}^k$ ($k \geq 1$). \omterm{def:b2:curves:length}{دالة طول القوس}
$s$ تقابل تفاضلي من الصنف $\mathcal{C}^k$ من $I$ على فترة $J$، والقوس
$\tilde\gamma =
\gamma \circ s^{-1}$ يحقق $\norm{\tilde\gamma'} = 1$ في كل نقطة. وإلى غاية انسحاب للوسيط وتوجيه، يكون
هذا التوسيم \emph{\omterm{def:b2:curves:length}{بطول القوس}} (أو \emph{ذو السرعة الواحدية}) وحيدا.
\end{theorem}

\begin{proof}
$s'(t) = \norm{\gamma'(t)} > 0$ والمقدار $s'$ من الصنف $\mathcal{C}^{k-1}$ (فهو مركب التطبيق $\gamma'$ من
الصنف $\mathcal{C}^{k-1}$ مع \omterm{def:b2:nvs:norm}{المعيار}، وهو أملس بعيدا عن $0$)، ومنه فإن $s$
من الصنف $\mathcal{C}^k$، ومتزايدة تماما، وتقابل على $J = s(I)$، وعكسها من الصنف
$\mathcal{C}^k$ بمبرهنة الدالة العكسية في البعد $1$ (مجلد السنة الأولى).
ثم
\[
\tilde\gamma'(\sigma)
= \frac{\gamma'(t)}{s'(t)}
= \frac{\gamma'(t)}{\norm{\gamma'(t)}},
\qquad t = s^{-1}(\sigma),
\]
وهي متجهة واحدية. وإذا كان $\hat\gamma = \gamma\circ\theta$ توسيما آخر ذا سرعة واحدية، فإن
$\abs{\theta'} = 1$، ومنه $\theta'
= \pm 1$ ثابت (بالاتصال)، أي $\theta(u) = \pm u + c$.
\end{proof}

\begin{remark}
\omterm{def:b2:curves:length}{طول القوس} هو الوسيط الذي يفصل الهندسة عن الديناميكا. فالمسار يمكن
قطعه بأي منحى سرعة --- وهو فيزياء الحركة --- لكن لكل سؤال ثابت عند
\omterm{def:b2:curves:reparam}{تغيير الوسيط} (الشكل، الانثناء، التقبيل) ساعة قانونية، هي المسافة
المقطوعة. ولهذا فإن جميع صيغ \omterm{thm:b2:curves:frenet2d}{الانحناء} أدناه \emph{تُعرَّف} بالسرعة
الواحدية ثم \emph{تُترجَم} إلى توسيمات كيفما كانت بقاعدة السلسلة:
فعوامل الترجمة قوى للمقدار $v = s'$، وتتبعها تتبعا صحيحا هو كل مضمون
\cref{prop:b2:curves:kappaformula}.
\end{remark}

\begin{example}[الدائرة واللولب]\label{ex:b2:curves:helix}
من أجل الدائرة $\gamma(t) = (R\cos t, R\sin t)$، $\norm{\gamma'} =
R$، ومنه $s = Rt$، وطول دورة كاملة هو
$2\pi R$. ومن أجل اللولب $\gamma(t) = (a\cos t,\ a\sin t,\ bt)$ مع $a > 0$، يكون $\norm{\gamma'(t)} = \sqrt{a^2 + b^2}$ ثابتا: فاللولب
يُقطع بسرعة ثابتة، و$s = t\sqrt{a^2 + b^2}$.
\end{example}

\begin{example}[طول القوس بالإحداثيات القطبية]\label{ex:b2:curves:polarlength}
المنحنى القطبي $r = r(\theta)$ هو القوس $\gamma(\theta) =
(r\cos\theta,\ r\sin\theta)$، مع
\[
\gamma'(\theta) = (r'\cos\theta - r\sin\theta,\
r'\sin\theta + r\cos\theta),
\qquad
\norm{\gamma'(\theta)}^2 = r'^2 + r^2
\]
(فالحدود المختلطة تتلاشى): ومنه فإن عنصر الطول القطبي هو $\sqrt{r^2 + r'^2}\,\dd\theta$.
ومن أجل القلبية $r = 1 +
\cos\theta$:
\[
r^2 + r'^2 = (1 + \cos\theta)^2 + \sin^2\theta = 2 +
2\cos\theta = 4\cos^2\tfrac\theta2,
\]
وعلى $\intcc{-\pi}{\pi}$ يكون جيب تمام نصف الزاوية موجبا، ومنه
\[
L = \int_{-\pi}^{\pi}2\cos\tfrac\theta2\,\dd\theta
= \Bigl[4\sin\tfrac\theta2\Bigr]_{-\pi}^{\pi}
= 4 - (-4) = 8 :
\]
وكما في قوس الدويري في \cref{exo:b2:curves:1}، فإن منحنى مبنيا من دوائر طوله ناطق،
دون أي $\pi$ في أي موضع. وتفكيك نصف الزاوية هو الحيلة المعتادة
لأطوال المنحنيات \omterm{def:b2:structures:generated}{المولَّدة} بالدوائر؛ وعندما يفشل (كما في الأهليلج)،
يكون الطول دالة جديدة حقا --- تكاملا إهليلجيا، خارج الصيغ المغلقة
الأولية.
\end{example}

\section{الانحناء في المستوي}

في كل هذا القسم، الأقواس من الصنف $\mathcal{C}^2$ ونظامية في المستوي
الإقليدي الموجه. ونوسّم \omterm{def:b2:curves:length}{بطول القوس} ونكتب $T(s) = \tilde\gamma'(s)$ \omterm{def:b2:curves:arc}{للمماس} الواحدي،
و$N(s)$ للمتجهة الواحدية المتعامدة مباشرة مع $T(s)$ (أي دوران $T$
بالزاوية $+\pi/2$).

\begin{theorem}[صيغ فرينيه في المستوي]\label{thm:b2:curves:frenet2d}
ليكن $\tilde\gamma$ قوسا ذا سرعة واحدية من الصنف $\mathcal{C}^2$ في المستوي الموجه.
توجد دالة متصلة $\kappa$، هي \emph{الانحناء}\index{انحناء} (الجبري)،
تحقق
\[
T'(s) = \kappa(s)\, N(s),
\qquad
N'(s) = -\kappa(s)\, T(s).
\]
\end{theorem}

\begin{proof}
بما أن $\norm{T(s)}^2 = 1$ من أجل كل $s$، فإن مفاضلة الجداء السلمي تعطي
$2\langle T'(s), T(s)\rangle = 0$: أي إن $T'(s)$ متعامد مع $T(s)$، ومنه فهو خطي مع $N(s)$
(فالبعد $2$)؛ ونكتب $T'(s)
= \kappa(s) N(s)$ مع $\kappa(s) = \langle T'(s), N(s)\rangle$ متصلة. وبالمثل $N' \perp N$،
ومنه $N' = \lambda T$؛ ومفاضلة $\langle T, N\rangle = 0$ تعطي $\langle T', N\rangle +
\langle T, N'\rangle = \kappa + \lambda = 0$.
\end{proof}

\begin{definition}\label{def:b2:curves:curvature}
عندما يكون $\kappa(s) \neq 0$، يكون \emph{نصف قطر الانحناء}\index{نصف قطر الانحناء}
هو $R(s) =
1/\abs{\kappa(s)}$ ويكون \emph{مركز الانحناء}\index{مركز الانحناء} هو
$\tilde\gamma(s) + \frac{1}{\kappa(s)} N(s)$؛ والدائرة ذات هذا المركز وهذا النصف قطر $R(s)$ هي
\emph{دائرة التقبيل}\index{دائرة التقبيل}، أي أفضل تقريب دائري
للمنحنى في $\tilde\gamma(s)$.
\end{definition}

\begin{example}[دائرة التقبيل للدالة الأسية]\label{ex:b2:curves:oscexp}
من أجل $y = \eu^x$ في النقطة $(0, 1)$: $f'(0) = f''(0) = 1$، ومنه، بصيغة التمثيل
البياني أدناه،
\[
\kappa(0) = \frac{1}{(1 + 1)^{3/2}} = \frac1{2\sqrt2},
\qquad R = 2\sqrt2 .
\]
\omterm{def:b2:curves:arc}{والمماس} الواحدي هو $T = \frac{(1, 1)}{\sqrt2}$، والناظم المباشر $N = \frac{(-1, 1)}{\sqrt2}$، \omterm{def:b2:curves:curvature}{ومركز الانحناء}
هو
\[
(0, 1) + 2\sqrt2\cdot\frac{(-1, 1)}{\sqrt2} = (-2,\ 3) :
\]
ومنه فإن معادلة \omterm{def:b2:curves:curvature}{دائرة التقبيل} هي $(x + 2)^2 + (y - 3)^2 =
8$. وللتحقق من دعوى ``أفضل
تقريب دائري'': فحل معادلة الدائرة بدلالة $y$ بجوار $(0,1)$ ونشره
يعطي $y = 1 + x + \frac{x^2}2 + O(x^3)$ --- وهو بالضبط نشر تايلور من الرتبة الثانية للدالة
$\eu^x$. \omterm{def:b2:curves:curvature}{فدائرة التقبيل} توافق القيمة والميل \emph{والمشتقة الثانية}؛
أما دائرة مماسة عادية فتوافق الأولين فقط.
\end{example}

\begin{example}[مطوّرة الدائرة مركزها]\label{ex:b2:curves:circleevolute}
من أجل الدائرة ذات نصف القطر $R$ المقطوعة عكس عقارب الساعة، لدينا
$\kappa = 1/R$ ويشير $N$ نحو المركز، ومنه فإن \omterm{def:b2:curves:curvature}{مركز الانحناء} $\tilde\gamma + \frac1\kappa N$ هو
مركز الدائرة، من أجل كل $s$: \omterm{def:b2:curves:curvature}{فدائرة التقبيل} للدائرة هي الدائرة
نفسها، وينهار محل مراكز \omterm{thm:b2:curves:frenet2d}{الانحناء} إلى نقطة. وتعاير هذه الحالة المنحلة
\cref{exo:b2:curves:6}: فهناك تكون سرعة المطوّرة $-\frac{\kappa'}{\kappa^2}N$، وهي تنعدم تطابقا إذا وفقط
إذا كان $\kappa$ ثابتا.
\end{example}

\begin{proposition}[الانحناء في توسيم كيفما كان]\label{prop:b2:curves:kappaformula}
من أجل قوس مستو نظامي من الصنف $\mathcal{C}^2$، $\gamma(t) = (x(t), y(t))$،
\[
\kappa(t)
= \frac{x'(t)\,y''(t) - y'(t)\,x''(t)}
       {\bigl(x'(t)^2 + y'(t)^2\bigr)^{3/2}} ,
\]
وبوجه خاص $\kappa = \dfrac{y''}{(1 + y'^2)^{3/2}}$ من أجل التمثيل البياني $y = f(x)$.
\end{proposition}

\begin{proof}
نكتب $v(t) = \norm{\gamma'(t)} = s'(t)$، ومنه $\gamma' = vT$ (بتركيب معطيات السرعة الواحدية مع $s$).
وبالمفاضلة،
\[
\gamma'' = v'T + v\,T'\cdot s' = v'T + v^2\kappa N .
\]
ولنأخذ الآن \omterm{def:b2:linalg:det}{محدد} $(\gamma', \gamma'')$ (في الأساس القانوني الموجه): بما أن $\det(T, T) = 0$
و$\det(T, N) = 1$،
\[
\det(\gamma', \gamma'') = \det(vT,\ v'T + v^2\kappa N)
= v^3\kappa .
\]
والطرف الأيسر هو $x'y'' - y'x''$، و$v^3 = (x'^2 +
y'^2)^{3/2}$. وحالة التمثيل البياني هي التوسيم
$t \mapsto (t,
f(t))$.
\end{proof}

\begin{example}[الدائرة والمستقيم والقطع المكافئ]\label{ex:b2:curves:kappaexamples}
للمستقيم $\kappa = 0$ (والعكس صحيح: فالشرط $T' = 0$ يعني أن $T$ ثابت،
ومنه $\tilde\gamma(s) = \tilde\gamma(0) + sT$، وهو مستقيم). وللدائرة ذات نصف القطر $R$ المقطوعة
عكس عقارب الساعة $\kappa = 1/R$: فمع $\gamma(t) = (R\cos t, R\sin t)$، تعطي الصيغة $\kappa =
R^2/R^3$. ومن أجل
القطع المكافئ $y = x^2/2$: $\kappa(x) = 1/(1 +
x^2)^{3/2}$، وهو أعظمي عند الرأس --- فالقطع
المكافئ أشد انثناء حيث يستدير.
\end{example}

\begin{remark}[مزالق شائعة حول الانحناء]
(1) يتغير \omterm{thm:b2:curves:frenet2d}{الانحناء} الجبري لقوس مستو إشارةً عند عكس توجيه القوس أو
توجيه المستوي: فوحدهما $\abs\kappa$ و$R = 1/\abs\kappa$ هندسيان محضان. والدائرة
المقطوعة مع عقارب الساعة انحناؤها $\kappa = -1/R$. (2) صيغة التمثيل البياني
$\kappa = f''/(1 + f'^2)^{3/2}$ تختار بصمت التوسيم بالمتغير $x$؛ وتطبيقها على منحن ليس
تمثيلا بيانيا بجوار النقطة (\omterm{def:b2:curves:arc}{بمماس} شاقولي) هو الغلط الكلاسيكي. (3)
في نقطة يكون فيها $\gamma' = 0$ لا يُعرَّف شيء --- لا $T$ ولا $\kappa$
--- وقد ينكسر المسار فعلا (\cref{ex:b2:curves:cusp})؛ فتحقق دائما من النظامية قبل
مفاضلة \omterm{def:b2:curves:arc}{المماس} الواحدي. (4) في الفضاء، يكون $\kappa =
\norm{T'} \geq 0$ بالاصطلاح: فلا
إشارة يمكن أن تخطئ فيها، لكن لا إشارة تستغلها كذلك --- فالمعلومات
من نوع الانعطاف تنتقل إلى \omterm{thm:b2:curves:frenet3d}{الالتواء}. (5) وأخيرا، فإن $\kappa$ مشتقة
\emph{بالنسبة إلى \omterm{def:b2:curves:length}{طول القوس}}: ففي حالة توسيم ليست سرعته واحدية، يكون
نسيان العامل $v^3$ في \cref{prop:b2:curves:kappaformula} أكثر الأخطاء تواترا عمليا.
\end{remark}

\begin{figure}[ht]
\centering
\begin{tikzpicture}[scale=1.6]
  % parabola y = x^2/2
  \draw[->] (-2.2,0) -- (2.4,0) node[below] {$x$};
  \draw[->] (0,-0.4) -- (0,2.4) node[left] {$y$};
  \draw[thick, ocDarkBlue, domain=-2.1:2.1, smooth, samples=60]
    plot (\x, {0.5*\x*\x});
  % osculating circle at vertex: kappa = 1 => R = 1, center (0,1)
  \draw[ocRed, dashed] (0,1) circle (1);
  \fill[ocRed] (0,1) circle (0.03) node[right, font=\small]
    {مركز الانحناء};
  \fill (0,0) circle (0.03);
  % tangent and normal at vertex
  \draw[->, thick] (0,0) -- (0.7,0) node[below, font=\small] {$T$};
  \draw[->, thick] (0,0) -- (0,0.7) node[left, font=\small] {$N$};
  % a second point
  \fill (1.2,0.72) circle (0.03);
  \draw[->, thick] (1.2,0.72) -- ++(0.448,0.5376)
    node[right, font=\small] {$T$};
  \draw[->, thick] (1.2,0.72) -- ++(-0.5376,0.448)
    node[above, font=\small] {$N$};
\end{tikzpicture}
\caption{القطع المكافئ $y = x^2/2$، ومعلم فرينيه المتحرك $(T,
N)$، \omterm{def:b2:curves:curvature}{ودائرة
التقبيل} عند الرأس (نصف قطرها $1$، لأن $\kappa(0) = 1$). ويدور المعلم مع
تحرك النقطة؛ \omterm{thm:b2:curves:frenet2d}{والانحناء} هو معدل ذلك الدوران لكل وحدة من \omterm{def:b2:curves:length}{طول القوس}.}
\label{fig:b2:curves:osculating}
\end{figure}

\begin{theorem}[الانحناء يحدد المنحنى]\label{thm:b2:curves:fundamental}
لتكن $\kappa \colon J \to \R$ متصلة. يوجد قوس ذو سرعة واحدية من الصنف $\mathcal{C}^2$ في
المستوي انحناؤه $\kappa$، وهو وحيد إلى غاية تقايس مباشر (دوران متبوع
بانسحاب).
\end{theorem}

\begin{proof}
\emph{الوجود.} نثبت $s_0 \in J$ ونضع $\varphi(s) =
\int_{s_0}^s \kappa(u)\,\dd u$، ثم
\[
\tilde\gamma(s) = \Bigl(\int_{s_0}^s \cos\varphi(u)\,\dd u,\
\int_{s_0}^s \sin\varphi(u)\,\dd u\Bigr).
\]
عندئذ يكون $T(s) = (\cos\varphi(s), \sin\varphi(s))$ متجهة واحدية، و$N(s) = (-\sin\varphi, \cos\varphi)$، و
\[
T'(s) = \varphi'(s)\,(-\sin\varphi, \cos\varphi)
= \kappa(s)\,N(s) :
\]
فالقوس ذو سرعة واحدية وانحناؤه $\kappa$.

\emph{الوحدانية.} ليكن $\gamma_1, \gamma_2$ قوسين ذوي سرعة واحدية لهما \omterm{thm:b2:curves:frenet2d}{الانحناء}
نفسه. يُرفع كل \omterm{def:b2:curves:arc}{مماس} واحدي إلى دالة زاوية، بإنشاء صريح: ننظر إلى
$T_j$ بوصفه العدد المركب $z_j = a_j + \iu b_j$ ذا الطويلة $1$، ونختار $\varphi_j(0)$
بحيث $z_j(0) = \eu^{\iu\varphi_j(0)}$، ونضع
\[
\varphi_j(s) = \varphi_j(0) + \int_0^s\det\bigl(T_j,
T_j'\bigr)(u)\,\dd u .
\]
وانطلاقا من $\abs{z_j} = 1$: $\operatorname{Re}(\conj{z_j}\,z_j') =
0$، ومنه $\conj{z_j}\,z_j' = \iu\det(T_j, T_j') =
\iu\,\varphi_j'$، أي $z_j' = \iu\varphi_j'z_j$؛ ثم
\[
\bigl(z_j\,\eu^{-\iu\varphi_j}\bigr)'
= \eu^{-\iu\varphi_j}\bigl(z_j' - \iu\varphi_j'z_j\bigr) = 0,
\]
ومنه $z_j = \eu^{\iu\varphi_j}$ في كل نقطة: أي $T_j =
(\cos\varphi_j, \sin\varphi_j)$ مع $\varphi_j$ من الصنف $\mathcal C^1$.
وعلاوة على ذلك $\det(T_j, T_j') = \det(T_j,
\kappa N_j) = \kappa$، ومنه $\varphi_j' = \kappa$. ومنه $\varphi_2 = \varphi_1 + c$ من أجل ثابت $c$:
أي إن $T_2$ هو $T_1$ بعد دورانه بالزاوية الثابتة $c$، ومنه
بالمكاملة، $\gamma_2 = \rho(\gamma_1) + w$ حيث $\rho$ هو الدوران بالزاوية $c$ و$w$
متجهة ثابتة.
\end{proof}

\begin{remark}
هذا هو النموذج الأولي ذو البعد الواحد \emph{لمبرهنة أساسية في
الهندسة}: فمجموعة كاملة من الثوابت المحلية (وهي هنا دالة واحدة)
تصنّف الكائن إلى غاية حركة صلبة. وتحتاج الصيغة ذات الأبعاد الثلاثة
أدناه إلى ثابتين.
\end{remark}

\begin{example}[الانحناء الثابت يعني الدائرة]\label{ex:b2:curves:constantkappa}
نأخذ $\kappa \equiv \kappa_0 > 0$ في صيغة الوجود: $\varphi(s) = \kappa_0 s$ و
\[
\tilde\gamma(s) = \Bigl(\frac{\sin\kappa_0s}{\kappa_0},\
\frac{1 - \cos\kappa_0s}{\kappa_0}\Bigr) :
\]
أي الدائرة ذات نصف القطر $1/\kappa_0$ ومركزها $(0,
1/\kappa_0)$، مقطوعة بسرعة
واحدية. وبحسب شق الوحدانية من المبرهنة، يكون \emph{كل} قوس ذي سرعة
واحدية وانحناء ثابت $\kappa_0$ قطعة من دائرة نصف قطرها $1/\kappa_0$ (أو
مستقيما إذا كان $\kappa_0 = 0$) --- وهو عكس الحساب في \cref{ex:b2:curves:kappaexamples}، والحالة
المستوية من \cref{exo:b2:curves:9}.
\end{example}

\begin{example}[إعادة بناء منحن انطلاقا من انحنائه]\label{ex:b2:curves:reconstruction}
أي منحن ذي سرعة واحدية نصف قطر انحنائه $R(s) = 1 +
s^2$؟ باتباع برهان الوجود
مع $\kappa(s) =
\frac1{1+s^2}$ و$s_0 = 0$: $\varphi(s) = \arctan s$، ومنه
\[
T(s) = (\cos\arctan s,\ \sin\arctan s)
= \Bigl(\frac{1}{\sqrt{1+s^2}},\
\frac{s}{\sqrt{1+s^2}}\Bigr),
\]
وبالمكاملة،
\[
\tilde\gamma(s) = \Bigl(\ln\bigl(s + \sqrt{1 + s^2}\bigr),\
\sqrt{1 + s^2} - 1\Bigr).
\]
وبوضع $x = \ln(s + \sqrt{1+s^2})$، أي $s = \sinh x$، تكون الإحداثية الثانية $\cosh x - 1$: فالمنحنى
هو \emph{السلسلية} $y = \cosh x - 1$. وهذا يغلق الحلقة مع \cref{exo:b2:curves:3}، حيث حسبنا
$R = \cosh^2 x = 1 +
\sinh^2 x = 1 + s^2$ مباشرة: فالمبرهنة الأساسية تضمن أن السلسلية هي المنحنى
\emph{الوحيد} بهذا المنحى من \omterm{thm:b2:curves:frenet2d}{الانحناء}، إلى غاية تقايس مباشر.
\end{example}

\begin{remark}[أين يُستعمل الانحناء فيما يلي]
التفكيك $\gamma'' = v'T + v^2\kappa N$ المحصل عليه في برهان \cref{prop:b2:curves:kappaformula} هو الحركية لكل حركة
منحنية: التسارع المماسي مقابل التسارع المركزي. ويعود \omterm{thm:b2:curves:frenet2d}{الانحناء} في
حالة السطوح (\cref{ch:b2:surfaces}) عبر انحناء المنحنيات المرسومة عليها، وحساب
المغلِّفات في مسألة نهاية الأسبوع لهذا الفصل --- المطوّرات والمنحنيات
الحارقة --- هو البصريات الهندسية لجبهات الموجة. ويستأنف مجلد السنة
الثالثة وجهة النظر الجوهرية من أجل المتنوعات الجزئية في $\R^n$.
\end{remark}

\section{معلم فرينيه في الفضاء}

ليكن الآن $\tilde\gamma \colon J \to \R^3$ قوسا ذا سرعة واحدية من الصنف $\mathcal{C}^3$ و\emph{ثنائي
النظامية}: أي $T'(s) \neq 0$ من أجل كل $s$. عندئذ يعرّف $\kappa(s) = \norm{T'(s)} > 0$
\emph{\omterm{thm:b2:curves:frenet2d}{الانحناء}} (ولا إشارة له في الفضاء: إذ لا توجيه ناظمي مفضل)،
ونضع:
\[
N(s) = \frac{T'(s)}{\kappa(s)}
\quad\text{(الناظم الرئيسي)},
\qquad
B(s) = T(s) \wedge N(s)
\quad\text{(الناظم الثنائي)} ,
\]
بحيث يكون $(T, N, B)$ معلما متعامدا مباشرا، هو \emph{معلم فرينيه}.
والمستوي المار بالنقطة $\tilde\gamma(s)$ والمولَّد بالمتجهتين $T, N$ هو
\emph{مستوي التقبيل}.

\begin{theorem}[صيغ فرينيه في الفضاء]\label{thm:b2:curves:frenet3d}
توجد دالة متصلة $\tau$، هي \emph{الالتواء}\index{التواء}، تحقق
\[
T' = \kappa N, \qquad
N' = -\kappa T + \tau B, \qquad
B' = -\tau N .
\]
\end{theorem}

\begin{proof}
الصيغة الأولى هي تعريف $N$. ولكل متجهة من المتجهات $T, N, B$ \omterm{def:b2:nvs:norm}{معيار}
ثابت $1$ وهي متعامدة مثنى مثنى؛ ومفاضلة العلاقات الست $\langle X, Y\rangle = \delta_{XY}$
تبين أن مصفوفة $(T', N', B')$ في الأساس $(T, N, B)$ متخالفة التناظر: فلدينا
فعلا $\langle X', Y\rangle + \langle X, Y'\rangle = 0$ و$\langle X', X\rangle = 0$. ومعاملها $(N, T)$ هو $\langle N', T\rangle =
-\langle N, T'\rangle = -\kappa$، ومعامل عمودها
$(T, B)$ هو $\langle T', B\rangle = \kappa\langle N, B\rangle = 0$. وتسمية المعامل الحر المتبقي $\tau = \langle N', B\rangle$ تعطي بالضبط
الصيغ الثلاث المكتوبة: فتخالف التناظر يملأ $\langle B', N\rangle = -\tau$ و$\langle B', T\rangle = 0$. واتصال
$\tau = \langle N', B\rangle$ واضح لأن $N'$ و$B$ متصلان (فالمقدار $\tilde\gamma$ من الصنف
$\mathcal{C}^3$، ومنه فإن $N = T'/\kappa$ من الصنف $\mathcal{C}^1$).
\end{proof}

\begin{example}[متجهة داربو]\label{ex:b2:curves:darboux}
تنضغط صيغ فرينيه الثلاث في صيغة واحدة. نضع $\omega(s) =
\tau\,T + \kappa\,B$ (وهي \emph{متجهة
داربو}). وباستعمال $B
\wedge T = N$ و$T \wedge N = B$ و$N \wedge B = T$:
\[
\omega \wedge T = \kappa\,N = T', \qquad
\omega \wedge N = \tau\,B - \kappa\,T = N', \qquad
\omega \wedge B = -\tau\,N = B' :
\]
فتتطور كل متجهة من المعلم وفق $X' = \omega \wedge X$، وهي البصمة الحركية
\emph{لدوران لحظي} متجهة سرعته الزاوية $\omega$. فالمعلم يدور بالمعدل
$\norm\omega = \sqrt{\kappa^2 + \tau^2}$ حول المحور المتحرك $\omega$؛ \omterm{thm:b2:curves:frenet2d}{والانحناء} هو مركبة الدوران على
الناظم الثنائي، \omterm{thm:b2:curves:frenet3d}{والالتواء} مركبته على \omterm{def:b2:curves:arc}{المماس}. وفي حالة اللولب، يكون
$\omega$ متجهة ثابتة موازية لمحور الأسطوانة --- وهذا بالضبط سبب
مبادرة معلم اللولب بانتظام. وتخالف تناظر مصفوفة فرينيه، المستغل في
\cref{exo:b2:curves:7}، هو الصيغة المصفوفية لهذه الحقيقة الهندسية الواحدة.
\end{example}

\begin{proposition}[الالتواء يقيس الاستواء]\label{prop:b2:curves:torsion}
يكون قوس ثنائي النظامية محتوى في مستو إذا وفقط إذا كان $\tau
\equiv 0$؛
وعندئذ يكون المستوي هو مستوي التقبيل (الثابت).
\end{proposition}

\begin{proof}
إذا كان $\tau \equiv 0$، فإن $B' = 0$، ومنه فإن $B$ متجهة واحدية ثابتة $B_0$،
و
\[
\frac{\dd}{\dd s}\langle \tilde\gamma(s), B_0\rangle
= \langle T(s), B_0\rangle = 0 :
\]
فالمقدار $\langle \tilde\gamma, B_0\rangle$ ثابت، ومنه يقع القوس في مستو متعامد مع $B_0$.
وبالعكس، إذا وقع القوس في مستو $P$، فإن $T$ و$T'$ (ومنه
$N$) موازيان لاتجاه $P$ من أجل كل $s$؛ ومنه فإن $B = T \wedge N$
أحد الناظمين الواحديين للمستوي $P$، وبما أنه \omterm{def:b2:metric:continuity}{متصل} فهو ثابت؛ ثم
إن $0 = B'
= -\tau N$ مع $N \neq 0$ يفرض $\tau \equiv 0$.
\end{proof}

\begin{example}[دائرة مائلة التواؤها معدوم]
القوس $\gamma(t) = \bigl(\cos t,\ \tfrac{\sin t}{\sqrt2},\
\tfrac{\sin t}{\sqrt2}\bigr)$ يقع في المستوي $y = z$، وهو دائرة الوحدة لذلك المستوي
(وللتحقق: $\norm{\gamma(t)} =
1$، ويُظهر الأساس المتعامد للمستوي $(1,0,0)$ و$(0, \tfrac1{\sqrt2}, \tfrac1{\sqrt2})$
التوسيم القياسي). ودون أي حساب فرينيه، تتنبأ \cref{prop:b2:curves:torsion} بأن $\tau \equiv 0$،
ويجب أن يكون الناظم الثنائي الثابت هو الناظم الواحدي للمستوي $\pm(0,
\tfrac1{\sqrt2}, -\tfrac1{\sqrt2})$.
\omterm{thm:b2:curves:frenet3d}{فالالتواء} لا يقيس ``الميل في الفضاء''؛ بل يقيس \emph{مغادرة} مستو.
ووحده المقدار $b$ غير المعدوم في اللولب أدناه ينتج التواء حقيقيا.
\end{example}

\begin{example}[اللولب]\label{ex:b2:curves:helixfrenet}
من أجل اللولب $\gamma(t) = (a\cos t, a\sin t, bt)$، $a > 0$، حسبنا $s = ct$ مع $c = \sqrt{a^2 + b^2}$. ثم
\[
T = \frac1c(-a\sin t,\ a\cos t,\ b),
\qquad
T' \cdot \frac{\dd t}{\dd s}
= \frac{1}{c^2}(-a\cos t, -a\sin t, 0),
\]
ومنه $\kappa = a/c^2 = a/(a^2 + b^2)$ و$N = (-\cos t, -\sin t,
0)$: فالناظم الرئيسي يشير أفقيا نحو المحور. ثم
$B = T \wedge N = \frac1c(b\sin t, -b\cos t, a)$، وتعطي $B'
\frac{\dd t}{\dd s} = \frac{b}{c^2}(\cos t, \sin t, 0) = -\tau N$
\[
\boxed{\ \kappa = \frac{a}{a^2 + b^2}, \qquad
\tau = \frac{b}{a^2 + b^2}. \ }
\]
والثابتان كلاهما ثابت --- ويمكن أن نبين، بالعكس، أن المنحنيات
الوحيدة ثنائية النظامية ذات $\kappa > 0$ الثابت و$\tau$ الثابت هي
اللوالب (والدوائر عندما $\tau = 0$). ولاحظ الإشارات: فالمقدار $b >
0$
يعطي لولبا أيمن اليد ذا التواء موجب.
\end{example}

\begin{example}[الانحناء والالتواء دون طول القوس؛ التكعيبي
الملتوي]\label{ex:b2:curves:twistedcubic}
إعادة التوسيم \omterm{def:b2:curves:length}{بطول القوس} مستحيلة عادة في صيغة مغلقة، ومنه يجب
استخراج الثوابت من المشتقات الخام. نكتب $v = \norm{\gamma'} = s'$؛ عندئذ $\gamma' =
vT$، وكما
في برهان \cref{prop:b2:curves:kappaformula}،
\[
\gamma'' = v'T + v^2\kappa N,
\qquad
\gamma' \wedge \gamma'' = v^3\kappa\,(T \wedge N) =
v^3\kappa\,B .
\]
وبأخذ المعايير ($\kappa \geq 0$ في الفضاء):
\[
\kappa = \frac{\norm{\gamma' \wedge \gamma''}}{v^3} .
\]
وبمفاضلة $\gamma''$ مرة أخرى وتحويل $N' =
v(-\kappa T + \tau B)$ (بقاعدة السلسلة عبر $s$)،
تأتي المركبة الوحيدة على $B$ من الحد الأخير:
\[
\gamma''' = \bigl(v'' - v^3\kappa^2\bigr)T +
\bigl(v'v\kappa + (v^2\kappa)'\bigr)N + v^3\kappa\tau\,B,
\]
ومنه، بالاقتران مع $\gamma' \wedge \gamma'' = v^3\kappa B$،
\[
\det(\gamma', \gamma'', \gamma''') = \langle\gamma' \wedge
\gamma'',\ \gamma'''\rangle = v^6\kappa^2\tau,
\qquad\text{أي}\qquad
\tau = \frac{\det(\gamma', \gamma'',
\gamma''')}{\norm{\gamma' \wedge \gamma''}^2} .
\]
وتطبيقا على \emph{التكعيبي الملتوي} $\gamma(t) = (t,\
t^2,\ t^3)$ عند $t = 0$: $\gamma' = (1, 0, 0)$،
$\gamma'' =
(0, 2, 0)$، $\gamma''' = (0, 0, 6)$، ومنه $v = 1$،
\[
\gamma' \wedge \gamma'' = (0, 0, 2), \qquad \kappa(0) = 2,
\qquad
\det(\gamma', \gamma'', \gamma''') = 12, \qquad \tau(0) =
\frac{12}{4} = 3 .
\]
والخلاصة: كلتا الصيغتين نسبة تُختصر فيها السرعة $v$ بالدرجة
المطلوبة بالضبط --- فالمقدار $\kappa$ يتحاكى كمشتقة ثانية لكل وحدة
طول، والمقدار $\tau$ كحجم مختلط لثلاث مشتقات لكل مساحة مربعة ---
وهذا هو \emph{سبب} كونهما هندسيين بينما $\gamma''$ نفسه ليس كذلك.
\end{example}

\begin{remark}[المبرهنة الأساسية للمنحنيات الفضائية]
كما في المستوي، يحدد الزوج $(\kappa, \tau)$ مع $\kappa > 0$ قوسا ثنائي النظامية
إلى غاية تقايس مباشر للفضاء $\R^3$: فصيغ فرينيه تشكل جملة تفاضلية
خطية للمعلم $(T, N, B)$، تنطبق عليها نظرية كوشي--ليبشيتز في \cref{ch:b2:diffeq}؛
ويُحافَظ على تعامد معلم الحل لأن مصفوفة المعاملات متخالفة التناظر
(بحجة مصفوفة غرام نفسها الواردة في \cref{exo:b2:curves:7})، ويُستعاد المنحنى
بمكاملة $T$. ونترك التفاصيل للقارئ تمرينا جوهريا لكنه مفيد.
\end{remark}

\section{الدراسة المحلية: الوضع بالنسبة إلى المماس}

\begin{proposition}[الشكل المحلي في نقطة نظامية]\label{prop:b2:curves:local}
ليكن $\gamma$ قوسا مستويا من الصنف $\mathcal{C}^k$ في $t_0$، وليكن $p$
أصغر دليل يحقق $\gamma^{(p)}(t_0) \neq 0$، وليكن $q$ أصغر دليل $> p$ بحيث لا يكون
$\gamma^{(q)}(t_0)$ خطيا مع $\gamma^{(p)}(t_0)$ (بافتراض وجودهما، $q \leq k$). وفي الأساس $(u, v) = (\gamma^{(p)}(t_0), \gamma^{(q)}(t_0))$
المرتكز في $\gamma(t_0)$، تعطي مبرهنة تايلور--يونغ الإحداثيات
\[
X(t) \sim \frac{(t - t_0)^p}{p!},
\qquad
Y(t) \sim \frac{(t - t_0)^q}{q!} .
\]
ولا تتعلق الصورة المحلية إلا بزوجية $p$ و$q$:
\begin{center}
\begin{tabular}{c c l}
$p$ فردي، $q$ زوجي & \emph{نقطة عادية} & لا يعبر المنحنى أي مستقيم،
ويبقى في جهة واحدة من \omterm{def:b2:curves:arc}{المماس}\\
$p$ فردي، $q$ فردي & \emph{نقطة انعطاف} & يعبر المنحنى مماسه\\
$p$ زوجي، $q$ فردي & \emph{نقطة رجوع من النوع الأول} & الفرعان
في جهتين متقابلتين من \omterm{def:b2:curves:arc}{المماس}\\
$p$ زوجي، $q$ زوجي & \emph{نقطة رجوع من النوع الثاني} & الفرعان
في الجهة نفسها\\
\end{tabular}
\end{center}
\end{proposition}

\begin{proof}
مبرهنة تايلور--يونغ من الرتبة $q$ (فالدالة $\gamma$ من الصنف
$\mathcal{C}^q$ بجوار $t_0$):
\[
\gamma(t) - \gamma(t_0)
= \sum_{j=p}^{q} \frac{(t-t_0)^j}{j!}\,\gamma^{(j)}(t_0)
+ o\bigl((t-t_0)^q\bigr).
\]
وبحكم اختيار $p$ و$q$، يكون كل $\gamma^{(j)}(t_0)$ بحيث $p
\leq j < q$ خطيا مع
$u$؛ وبجمع المركبات في الأساس $(u, v)$: $X(t) = \frac{(t-t_0)^p}{p!}(1 + o(1))$ و$Y(t) = \frac{(t-t_0)^q}{q!}(1 + o(1))$. وجدول
إشارتي $X$ و$Y$ من أجل $t \gtrless t_0$ --- الذي تحكمه الزوجية بالضبط
--- يعطي الصور الأربع: فمثلا إذا كان $p$ زوجيا، كان $X > 0$ في
الجهتين معا (فالفرعان يغادران في الاتجاه $+u$: أي نقطة رجوع)،
وتنقلب جهة \omterm{def:b2:curves:arc}{المماس} ($\operatorname{sign} Y$) عندما يكون $q$ فرديا.
\end{proof}

\begin{example}
من أجل $\gamma(t) = (t^2, t^3)$ في $t_0 = 0$ (\cref{ex:b2:curves:cusp}): $\gamma'' (0)= (2, 0)$، $\gamma'''(0) =
(0, 6)$، ومنه $p = 2$،
$q = 3$: أي نقطة رجوع من النوع الأول، وهي الصورة المألوفة للقطع
المكافئ نصف التكعيبي. ومن أجل $\gamma(t) = (t,
t^3)$ في $0$: $p = 1$، $q = 3$:
أي انعطاف --- فالتكعيبي يعبر مماسه.
\end{example}

\begin{remark}[آفاق داخل هذا المجلد]
تغذي المنحنيات الفصول التالية بثلاث طرائق. فالمنحنيات المرسومة على
سطح تعرّف مستوياته المماسة وصورته الأساسية الأولى (\cref{ch:b2:surfaces})، وتُحسب
أطوالها بقصر المتري المحيط --- فالفصل الآتي هو إلى حد بعيد هذا
الفصل بعد نسبته إلى سطح. وحساب المغلِّفات في مسألة نهاية الأسبوع
يلتقي بالتكاملات المزدوجة في \cref{ch:b2:multint}، حيث تُعاد حسبة مساحة \omterm{pb:b2:curves:1}{النجمية}
بصيغة غرين (\cref{exo:b2:multint:5}) --- منحنى واحد، ونظريتان، وجوابان متطابقان.
وقد استعملت جملة فرينيه أصلا المعادلات التفاضلية الخطية في \cref{ch:b2:diffeq}
(الوجود والوحدانية وحجة الحفاظ على التعامد في \cref{exo:b2:curves:7}): فالمبرهنة
الأساسية للمنحنيات مبرهنة معادلات تفاضلية في ثوب هندسي.
\end{remark}

\begin{remark}[طريقة: إجراء الدراسة المحلية]
عمليا يكون التصنيف روتينا من أربع خطوات. \emph{أولا}، فاضل في $t_0$
حتى تظهر أول مشتقة غير معدومة: فدليلها $p$ وقيمتها المتجهة $u$.
\emph{ثانيا}، تابع المفاضلة حتى تظهر مشتقة غير خطية مع $u$:
فدليلها $q$ ومتجهتها $v$. \emph{ثالثا}، اقرأ الزوجيتين $(p, q)$
في الجدول. \emph{رابعا}، ارسم: فالمنحنى يغادر في الاتجاه $+u$ إذا
كان $p$ فرديا (وفي الاتجاه $u$ ثم عائدا في الاتجاه $u$
إذا كان $p$ زوجيا)، في الجهة من $v$ التي تمليها إشارة $Y$.
وتحذيران. المعلم $(u, v)$ \emph{ليس} متعامدا عادة --- فالجدول يصف
المواضع بالنسبة إلى \omterm{def:b2:curves:arc}{المماس}، لا الزوايا ولا المسافات، فلا تقرأ
\omterm{thm:b2:curves:frenet2d}{الانحناء} من الصورة. والمشتقات الوسيطة الخطية مع $u$ مسموح بها
بين الرتبتين $p$ و$q$ (فهي لا تفعل إلا إزاحة نشر $X$)؛ أما
ما \emph{يجب ألا} يحدث فهو التوقف عند أول مشتقة غير معدومة وتخمين
$q = p + 1$: فمن أجل $\gamma(t)
= (t^2, t^4 + t^5)$ يكون التخمين الساذج $q = 3$ خاطئا، لأن $\gamma^{(3)}(0)$
لا يزال خطيا مع $\gamma''(0)$ --- وهذا بالضبط هو \cref{exo:b2:curves:5}.
\end{remark}

\section{تمارين}

\begin{exercise}[$\star$]\label{exo:b2:curves:1}
احسب طول قوس واحد من الدويري $\gamma(t) = (t -
\sin t,\ 1 - \cos t)$، $t \in [0, 2\pi]$. \emph{(استعمل $1 - \cos t =
2\sin^2(t/2)$.)}
\end{exercise}

\begin{exercise}[$\star$]\label{exo:b2:curves:2}
احسب انحناء الأهليلج $\gamma(t) = (a\cos t,\ b\sin
t)$ ($a > b > 0$) وحدد نقط \omterm{thm:b2:curves:frenet2d}{الانحناء} الأعظمي
والأصغري.
\end{exercise}

\begin{exercise}[$\star$]\label{exo:b2:curves:3}
بيّن أن طول قوس التمثيل البياني للدالة $f(x) = \cosh x$ على $[0, x]$ يساوي
$\sinh x$، واحسب انحناء هذا المنحنى (وهو \emph{السلسلية}). وتحقق من أن
$R(x) = 1/\kappa(x) =
\cosh^2 x$.
\end{exercise}

\begin{exercise}[$\star\star$]\label{exo:b2:curves:4}
(الحلزون اللوغاريتمي) لتكن $\gamma(t) = e^{t}(\cos t,\ \sin t)$، $t
\in \R$.
بيّن أن الزاوية بين $\gamma(t)$ و$\gamma'(t)$ ثابتة، واحسب طول قوس $\gamma$ على
$(-\infty, 0]$ (وهو منته!)، ثم احسب \omterm{thm:b2:curves:frenet2d}{الانحناء}.
\end{exercise}

\begin{exercise}[$\star\star$]\label{exo:b2:curves:5}
عيّن $p$ و$q$ والشكل المحلي (عادي، انعطاف، نقطة رجوع) للقوس
$\gamma(t) = (t^2,\ t^4 + t^5)$ في $t = 0$، وللقوس $\gamma(t) = (t^3,\ t^4)$ في $t = 0$.
\end{exercise}

\begin{exercise}[$\star\star$]\label{exo:b2:curves:6}
ليكن $\gamma$ قوسا مستويا ذا سرعة واحدية يحقق $\kappa(s) > 0$ من أجل كل
$s$، وليكن $c(s) = \gamma(s) + \frac{1}{\kappa(s)}N(s)$ \omterm{def:b2:curves:curvature}{مركز الانحناء} (والمنحنى $c$ هو
\emph{المطوّرة}). بافتراض أن $\kappa$ من الصنف $\mathcal{C}^1$، بيّن أن $c'(s) =
-\frac{\kappa'(s)}{\kappa(s)^2}N(s)$:
أي إن المطوّرة مماسة لنواظم $\gamma$.
\end{exercise}

\begin{exercise}[$\star\star\star$]\label{exo:b2:curves:7}
لتكن $A(s)$ عائلة متصلة من المصفوفات المتخالفة التناظر من الحجم
$3 \times 3$، ولتكن $F' = F A$ حلا مصفوفيا مع $F(s_0)$ متعامدة. بيّن أن $F(s)$
متعامدة من أجل كل $s$. \emph{(فاضل $G = F F^{\mathsf T}$ واستعمل الوحدانية في
مبرهنة كوشي--ليبشيتز.)} واشرح صلة ذلك بجملة فرينيه.
\end{exercise}

\begin{exercise}[$\star\star\star$]\label{exo:b2:curves:8}
(\omterm{thm:b2:curves:frenet2d}{الانحناء} الكلي لمنحن مغلق محدب) ليكن $\tilde\gamma$ قوسا مستويا مغلقا ذا
سرعة واحدية من الصنف $\mathcal{C}^2$ وطوله $L$ (بحيث $\tilde\gamma(s + L) = \tilde\gamma(s)$)، مقطوعا
مرة واحدة عكس عقارب الساعة. باستعمال دالة الزاوية $\varphi$ التي تحقق
$T =
(\cos\varphi, \sin\varphi)$ من \cref{thm:b2:curves:fundamental}، اشرح لماذا يكون $\varphi(L) -
\varphi(0)$ مضاعفا للمقدار $2\pi$،
وبيّن أن $\int_0^L \kappa(s)\,\dd s = \varphi(L) - \varphi(0)$. (ومن أجل دائرة نصف قطرها $R$: $\int \kappa = \frac1R \cdot 2\pi R = 2\pi$. وتؤكد
مبرهنة المماسات الدوارة القيمة $2\pi$ من أجل كل منحن مغلق بسيط؛
ولا يُطلب منك البرهان عليها.)
\end{exercise}

\begin{exercise}[$\star\star\star$]\label{exo:b2:curves:9}
بيّن أن منحنى فضائيا ثنائي النظامية ذا $\kappa > 0$ ثابت و$\tau = 0$ هو (قوس
من) دائرة نصف قطرها $1/\kappa$. \emph{(استعمل \cref{prop:b2:curves:torsion}، ثم بيّن أن المركز
$\gamma + \frac1\kappa N$ ثابت.)}
\end{exercise}

\begin{exercise}[$\star$]\label{exo:b2:curves:10}
احسب طول قوس القطع المكافئ $y = x^2/2$ على $\intcc0a$ وبيّن أنه يساوي
\[
\tfrac12\Bigl(a\sqrt{1 + a^2} + \ln\bigl(a + \sqrt{1 +
a^2}\bigr)\Bigr).
\]
\end{exercise}

\begin{exercise}[$\star\star$]\label{exo:b2:curves:11}
ليكن $\gamma$ قوسا نظاميا من الصنف $\mathcal C^2$ في $\R^n$ تمر جميع مماساته
بنقطة ثابتة $P$. برهن على أن مسار $\gamma$ محتوى في مستقيم.
\emph{(وسّم \omterm{def:b2:curves:length}{بطول القوس}، واكتب $\gamma(s) +
\lambda(s)T(s) = P$ وفاضل.)}
\end{exercise}

\begin{exercise}[$\star\star\star$]\label{exo:b2:curves:12}
(المبرهنة الأساسية للمنحنيات الفضائية) لتكن $\kappa > 0$ و$\tau$ دالتين
متصلتين على فترة $J$. نفّذ برنامج الملاحظة التالية للمبرهنة
\cref{ex:b2:curves:helixfrenet}: (1) بيّن أن الجملة الخطية $F' = FA(s)$، حيث $A(s)$ مصفوفة فرينيه
المتخالفة التناظر المبنية من $\kappa, \tau$ و$F(s_0)$ معلم متعامد مباشر، لها
حل شامل وحيد يبقى معلما متعامدا مباشرا؛ (2) أنشئ منحنى ذا سرعة
واحدية ثنائي النظامية انحناؤه $\kappa$ والتواؤه $\tau$؛ (3) برهن على
الوحدانية إلى غاية تقايس مباشر للفضاء $\R^3$.
\end{exercise}

\section{مسألة: المغلِّفات --- النجمية ومطوّرتان ومنحنى حارق}

\begin{omfigure}
\begin{tikzpicture}[scale=2.2]
  \draw[->] (-1.15,0) -- (1.25,0) node[below] {$x$};
  \draw[->] (0,-1.15) -- (0,1.25) node[left] {$y$};
  \draw[omDef] (1,0) -- (0,0.0) (0.966,0) -- (0,0.259)
    (0.866,0) -- (0,0.5) (0.707,0) -- (0,0.707)
    (0.5,0) -- (0,0.866) (0.259,0) -- (0,0.966);
  \draw[omThm, very thick]
    plot[domain=0:360, samples=120]
    ({cos(\x)*cos(\x)*cos(\x)}, {sin(\x)*sin(\x)*sin(\x)});
\end{tikzpicture}

{\small سلّم طوله $1$ ينزلق على جدار (بالأزرق مواضعه) ولا يعبر
أبدا \omterm{pb:b2:curves:1}{النجمية} $x^{2/3} + y^{2/3} = 1$ (بالأحمر): \omterm{pb:b2:curves:1}{فالنجمية} هي \emph{مغلِّف} عائلة
القطع، وهي مماسة لكل واحدة منها.}
\end{omfigure}

\begin{problem}[{مسألة نهاية الأسبوع --- آلة المغلِّفات وأربعة
منحنيات كلاسيكية}]\label{pb:b2:curves:1}
عادة ما تفشل عائلة مستقيمات ذات وسيط واحد في تغطية المستوي تغطية
منتظمة: فالمستقيمات تتكدس على منحن \omterm{def:b2:curves:arc}{مماس} لها جميعا، هو
\emph{مغلِّفها}\index{مغلف}. وتجعل أشعة الضوء المغلِّفات مرئية على
صورة \emph{منحنيات حارقة} --- وهو المنحنى المضيء ذو نقط الرجوع في
فنجان القهوة. وتبني هذه المسألة آلة المغلِّفات العامة، ثم تشغّلها
أربع مرات: السلم المنزلق (\omterm{pb:b2:curves:1}{النجمية})، ونواظم القطع المكافئ ونواظم
الدويري (المطوّرات، مع نوّاس هويغنز في النهاية)، والمنحنى الحارق في
فنجان القهوة (\omterm{pb:b2:curves:1}{الكلوية}). وفي كل ما يلي، يرمز $D_t$ إلى المستقيم ذي
المعادلة $a(t)\,x + b(t)\,y = c(t)$، حيث $a, b, c$ دوال من الصنف $\mathcal C^2$ تحقق $(a(t), b(t)) \neq (0,0)$،
و$\Delta(t) = a(t)b'(t) - a'(t)b(t)$.

\medskip
\noindent\textbf{الجزء الأول --- آلة المغلِّفات.}
\begin{enumerate}
  \item لنفترض $\Delta(t) \neq 0$. بيّن أن \emph{الجملة المميزة}
        \[
        \begin{cases} a(t)\,x + b(t)\,y = c(t)\\
        a'(t)\,x + b'(t)\,y = c'(t)\end{cases}
        \]
        لها حل وحيد $E(t) = (x(t), y(t))$، يعطى بالعلاقتين $x = \dfrac{cb' - c'b}{\Delta}$ و$y = \dfrac{ac' -
        a'c}{\Delta}$.
  \item لنفترض علاوة على ذلك أن $E$ من الصنف $\mathcal C^1$ بجوار
        $t$ مع $E'(t) \neq 0$. بمفاضلة المعادلة الأولى من الجملة، بيّن
        $a(t)\,x'(t) +
        b(t)\,y'(t) = 0$، واستنتج أن المنحنى $E$ يمر بنقطة من $D_t$
        \emph{وباتجاه} $D_t$: أي إن العائلة مماسة للمنحنى $E$،
        الذي يسمى \emph{\omterm{pb:b2:curves:1}{مغلِّفها}}.
  \item للتحقق: مماسات القطع المكافئ $y =
        x^2/2$ في النقط $(t, t^2/2)$ هي
        $tx - y =
        t^2/2$. تحقق من أن آلة المغلِّفات تعيد القطع المكافئ نفسه.
  \item (مغلِّف النواظم) ليكن $\gamma$ ذا سرعة واحدية مع $\kappa(s) \neq 0$.
        الناظم في $\gamma(s)$ هو $\{M : \langle M - \gamma(s), T(s)
        \rangle = 0\}$. بيّن أن جملته المميزة تفرض
        $\langle M - \gamma(s), N(s)\rangle =
        1/\kappa(s)$، ومنه أن النقطة المميزة هي \omterm{def:b2:curves:curvature}{مركز الانحناء}:
        \emph{فمغلِّف النواظم هو المطوّرة}، وهو ما يستعيد \cref{exo:b2:curves:6}.
        وتحقق من $\Delta(s) =
        \kappa(s)$.
  \item انحلالان. من أجل الحزمة $D_\theta : x\cos
        \theta + y\sin\theta = 0$، بيّن أن النقطة المميزة هي
        المبدأ من أجل كل $\theta$ (فينهار ``المغلِّف'' إلى نقطة، ولدينا
        $E' = 0$: فالسؤال 2 لا ينطبق). ومن أجل عائلة مستقيمات متوازية
        ($a, b$ ثابت)، بيّن $\Delta \equiv 0$ وأن الجملة المميزة غير متوافقة
        عموما: فلا مغلِّف.
\end{enumerate}

\medskip
\noindent\textbf{الجزء الثاني --- السلم المنزلق \omterm{pb:b2:curves:1}{والنجمية}.} قطعة
طولها $1$ تنزلق وأحد طرفيها $P_t =
(\cos t, 0)$ على الأرض والآخر $Q_t = (0, \sin t)$ على
الجدار، $t \in \intoo0{\pi/2}$.
\begin{enumerate}[resume]
  \item بيّن أن معادلة المستقيم $(P_tQ_t)$ هي $x\sin t +
        y\cos t = \sin t\cos t$، وأن آلة المغلِّفات
        تعطي النقطة المميزة
        \[
        E(t) = (\cos^3 t,\ \sin^3 t) :
        \]
        وهي \emph{النجمية}\index{نجمية}، ذات المعادلة الضمنية $x^{2/3} + y^{2/3} = 1$
        (ممتدة إلى الأرباع الأخرى بالتناظر).
  \item بيّن أن $E'(0) = 0$، واستعمل التصنيف المحلي (\cref{prop:b2:curves:local}) لتبين أن
        للنجمية نقطة رجوع من النوع الأول في $(1, 0)$ --- وكذلك في
        نقطها الأربع الواقعة على المحورين.
  \item احسب $\norm{E'(t)} = \tfrac32\abs{\sin 2t}$ واستنتج أن الطول الكلي للنجمية هو $6$.
  \item أين يلامس السلم \omterm{pb:b2:curves:1}{النجمية}؟ بيّن $E(t) =
        P_t + \sin^2 t\,(Q_t - P_t)$: فنقطة التماس تقسم
        السلم بالنسبة $\sin^2 t : \cos^2
        t$، وتجتاحه من طرف إلى طرف مع انزلاق
        السلم.
  \item احسب المساحة التي تحصرها \omterm{pb:b2:curves:1}{النجمية}: بيّن أن مساحة الربع الأول
        هي $3\int_0^{\pi/2}\sin^4
        t\cos^2 t\,\dd t$، واحسب التكامل بالتخطيط ($\sin^2 2t = \tfrac{1 - \cos 4t}2$)، واستنتج أن
        المساحة الكلية هي $3\pi/8$.
\end{enumerate}

\medskip
\noindent\textbf{الجزء الثالث --- مطوّرة القطع المكافئ.} ليكن
$\gamma(t) = (t, t^2/2)$.
\begin{enumerate}[resume]
  \item بيّن أن معادلة الناظم في $\gamma(t)$ هي $x + t\,y = t + t^3/2$.
  \item شغّل آلة المغلِّفات: بيّن أن مغلِّف النواظم هو
        \[
        E(t) = \Bigl(-t^3,\ 1 + \tfrac32 t^2\Bigr),
        \]
        ومعادلته الضمنية $x^2 = \tfrac8{27}(y - 1)^3$: أي قطع مكافئ نصف تكعيبي.
  \item تحقق متقاطعا مع السؤال 4: احسب \omterm{def:b2:curves:curvature}{مركز الانحناء} $\gamma(t) + \frac1{\kappa(t)}N(t)$ انطلاقا
        من
        $\kappa(t) = (1 + t^2)^{-3/2}$
        (\cref{ex:b2:curves:kappaexamples}) واستعد النقطة نفسها.
  \item بيّن أن للمطوّرة نقطة رجوع من النوع الأول في $(0, 1)$، وهو
        \omterm{def:b2:curves:curvature}{مركز الانحناء} عند الرأس --- أي النقطة التي يكون فيها $\kappa$
        أقصويا، كما تتنبأ به الصيغة $c' = -\frac{\kappa'}{\kappa^2}N$ في \cref{exo:b2:curves:6}.
  \item كم ناظما للقطع المكافئ يمر بنقطة معطاة $(x_0, y_0)$؟ بيّن أن
        الجواب تحكمه المعادلة التكعيبية $\tfrac{t^3}2 + (1 - y_0)\,t - x_0 = 0$؛ وعالج حالة المحور
        $x_0 = 0$ معالجة كاملة (ناظم واحد من أجل $y_0 < 1$، وثلاثة من أجل
        $y_0 > 1$)، وفسّر المطوّرة بوصفها منحنى الانتقال.
\end{enumerate}

\medskip
\noindent\textbf{الجزء الرابع --- المنحنى الحارق في فنجان القهوة.}
أشعة متوازية اتجاهها $(1, 0)$ تصطدم بالوجه الداخلي للمرآة الدائرية
$x^2 + y^2 = 1$؛ والشعاع الذي يصيب $P_\theta =
(\cos\theta, \sin\theta)$ ينعكس وفق قانون الانعكاس.
\begin{enumerate}[resume]
  \item انطلاقا من التناظر المرآتي حول الناظم (أي نصف القطر)، برر
        أن الاتجاه المنعكس هو $v = u -
        2\langle u, n\rangle n$ مع $u = (1,0)$ و$n =
        (\cos\theta, \sin\theta)$، ثم احسب
        $v =
        -(\cos2\theta, \sin2\theta)$.
  \item بيّن أن الشعاع المنعكس يقع على المستقيم
        \[
        x\sin 2\theta - y\cos 2\theta = \sin\theta .
        \]
  \item شغّل آلة المغلِّفات ($\Delta = 2$): بيّن أن المنحنى الحارق هو
        \[
        E(\theta) = \Bigl(\tfrac{3\cos\theta -
        \cos3\theta}4,\ \tfrac{3\sin\theta -
        \sin3\theta}4\Bigr),
        \]
        وهو \emph{الكلوية}\index{كلوية}.
  \item احسب $E'(\theta) = \tfrac32\sin\theta\,
        (\cos2\theta, \sin2\theta)$؛ وتحقق من أن اتجاه \omterm{def:b2:curves:arc}{المماس} هو اتجاه الشعاع
        المنعكس (السؤال 16)، وحدد نقطتي الرجوع $(\pm\tfrac12, 0)$، وبيّن أن
        الشعاع المنعكس يقطع المحور $y = 0$ في $x = \frac1{2\cos\theta}$ --- ومنه فإن
        الأشعة القريبة من المحور تتجمع في $x = \tfrac12$: أي إن البعد
        البؤري $R/2$ لمرآة نصف قطرها $R$.
  \item بيّن أن الطول الكلي للكلوية هو $6$، وأنه بجوار $\theta = 0$،
        \[
        E(\theta) - \bigl(\tfrac12, 0\bigr) =
        \bigl(\tfrac34\theta^2 + o(\theta^2),\ \theta^3 +
        o(\theta^3)\bigr) :
        \]
        أي نقطة رجوع من النوع الأول، متجهة على طول المحور.
  \item اشرح في فقرة واحدة لماذا يكون المنحنى الحارق مضيئا: فبكل
        نقطة تقع خارجه مباشرة يمر شعاعان منعكسان، وفي كل نقطة عليه
        تكون الأشعة ``مركّزة تركيزا لانهائيا'' (فالتطبيق
        $(\theta,
        \text{المسافة على طول الشعاع}) \mapsto \R^2$ له نقطة حرجة على
        المغلِّف بالضبط).
\end{enumerate}

\medskip
\noindent\textbf{الجزء الخامس --- هويغنز: الدويري مطوّرة نفسه.}
ليكن $\gamma(t) = (t - \sin t,\ 1 - \cos t)$، $t \in
\intoo0{2\pi}$، وهو قوس واحد من الدويري.
\begin{enumerate}[resume]
  \item احسب $\kappa(t) = -\dfrac1{4\sin(t/2)}$ \omterm{def:b2:curves:curvature}{ومركز الانحناء}؛ وبيّن أن المطوّرة هي
        \[
        c(t) = (t + \sin t,\ \cos t - 1),
        \]
        وأن التعويض $t = u + \pi$ يُظهرها الدويري الأصلي بعد انسحابه
        بالمتجهة $(\pi, -2)$: \emph{فمطوّرة الدويري دويري مطابق} (هويغنز).
  \item تحقق من أن \omterm{def:b2:curves:curvature}{نصف قطر الانحناء} عند القمة $t =
        \pi$ يساوي $4$،
        وهو نصف الطول $8$ لقوس واحد (\cref{exo:b2:curves:1})؛ وحدد نقطة رجوع
        المطوّرة أسفل القمة مباشرة، على المسافة $4$.
  \item (خاصية الخيط) ليكن $\gamma$ ذا سرعة واحدية مع $\kappa > 0$، و$\kappa$
        من الصنف $\mathcal C^1$ و$R = 1/\kappa$ رتيبة تماما. باستعمال $c' = R'N$، بيّن
        أن طول قوس المطوّرة بين $c(s_0)$ و$c(s_1)$ هو $\abs{R(s_1) - R(s_0)}$.
        وفسّر ذلك: فخيط مشدود يُفكّ عن المطوّرة، طوله $R(s_0)$ في
        البداية، يرسم طرفه الحر المنحنى الأصلي --- ومنه فإن نوّاسا
        يتأرجح بين خدّين دويريين من ساعة هويغنز يرسم دويريا.
  \item تركيب. أنتجت آلة الجزء الأول نجمية وقطعا مكافئا نصف تكعيبي
        وكلوية ودويريا. ومن أجل كل عائلة من العائلات الأربع، اذكر
        في جملة واحدة أين تحققت الفرضيتان $\Delta \neq 0$ و$E' \neq 0$ أو أين
        سقطتا، وأي حدث هندسي (نقطة رجوع، بؤرة، انحلال) دلّ عليه كل
        سقوط للفرضية $E' \neq 0$. وأين يجب أن تظهر القيم القصوى للانحناء
        على مغلِّف النواظم، ولماذا؟
\end{enumerate}
\end{problem}
